The Sirius star, located in the constellation of Canis Major, is the brightest
star in the night sky of the Earth. What the observer's eye sees as a single
star is actually a binary star system. The high brightness of Sirius is a
consequence of two facts: its intrinsic luminosity and its proximity to the
Earth.

The Mizar multiple star system, in the constellation of Ursa Major, consists
of 4 stars seen along the same line of sight from the Earth. Some of these
stars form a gravitationally bound system.

Let's assume that an observer (observer A) is located on one of the planets of
the Sirius system. \emph{Determine}:

\begin{description}
\item[a)] The magnitude of the Sun as seen by observer A
  ($m_{\mathrm{Sun,Planet}}$).
\item[b)] The magnitude of Sirius star system as seen by the observer A.
  ($m_{\mathrm{SY,Planet}}$)
\item[c)] The combined intrinsic luminosity of the Mizar system,
  $L_{\mathrm{Mizar}}$;
\item[d)] the average distance between gravitationally bound stars of the
  Mizar system and Earth, $\ell$
  % sic: the symbol in the original prints as a stray grave accent; it is
  % most plausibly a script ell.
\item[e)] The geocentric angular distance between Mizar system and Sirius,
  $\Delta\theta$;
\item[f)] The physical distance between the gravitationally bound stars of
  the Mizar system and the observer A. ($d_{\mathrm{Mizar,Planet}}$)
\item[g)] The magnitude of the entire Mizar system as seen by the observer A.
  ($m_{\mathrm{Mizar,Planet}}$)
\end{description}

Also estimate amount of errors in all your answers.

The following data may be used:
\begin{itemize}
\item $d_{\mathrm{Sirius,Earth}} = 2.6$\,pc --- the Sirius--Earth distance;
\item $m_{\mathrm{Sirius,Earth}} = -1.46^{\mathrm{m}}$ --- the apparent
  magnitude of Sirius measured from the Earth;
\item $d_{\mathrm{Sun,Earth}} = 1$\,AU --- the Sun--Earth distance;
\item $m_{\mathrm{Sun,Earth}} = -26.78^{\mathrm{m}}$ --- the apparent
  magnitude of the Sun as seen from Earth;
\item $d_{\mathrm{Sirius,Planet}} = 10$\,AU --- distance between Sirius and
  its planet where the observer A is located.
\end{itemize}

In the table below information for the stars from the Mizar system as
measured from the Earth is given.

\begin{center}
\begin{tabular}{cccc}
\hline
Star & Name of the & Apparent & Parallax \\
number & star & magnitude & (milli arc seconds) \\
\hline
1 & Alcor & $3.99 \pm 0.01$ & $39.91 \pm 0.13$ \\
2 & Mizar A & $2.23 \pm 0.01$ & $38.01 \pm 1.71$ \\
3 & Mizar B & $3.86 \pm 0.01$ & $38.01 \pm 1.71$ \\
4 & Sidus Ludoviciana & $7.56 \pm 0.01$ & $8 \pm 4$ \\
\hline
\end{tabular}
\end{center}

The equatorial coordinates of Mizar system ($\sigma_1$) and respectively of
Sirius ($\sigma_2$), located on the heliocentric map are:
\[ \alpha_{\mathrm{Mizar}} = \alpha_1 =
   13^{\mathrm{h}}23^{\mathrm{min}}55.5^{\mathrm{s}}; \quad
   \delta_{\mathrm{Mizar}} = \delta_1 = 54^\circ 55' 31''; \quad
   \alpha_{\mathrm{Sirius}} = \alpha_2 = 6^{\mathrm{h}}45^{\mathrm{min}};
   \quad
   \delta_{\mathrm{Sirius}} = \delta_2 = -16^\circ 43'. \]

Note
\[ \ln(1-x) \approx -x \text{ for } x \ll 1 \]
\[ e^x \approx 1 + x \text{ for } x \ll 1 \]
